\documentclass[
    language=english,
    doctype=thesis-proposal,
    institution=none,
]{../../omnilatex}

\RequirePackage{config/document-settings}
\addbibresource{../../bib/bibliography.bib}

\begin{document}

\maketitle

\section*{Abstract}
Bayesian optimisation is a sample-efficient approach to global optimisation of
expensive black-box functions.  However, its effectiveness degrades sharply in
dimensions beyond approximately 20, owing to the curse of dimensionality and
the difficulty of modelling complex posterior distributions.  This proposal
outlines a doctoral research plan to develop novel adaptive sampling strategies
that combine dimensionality reduction, local surrogate modelling, and
information-theoretic acquisition functions to enable scalable Bayesian
optimisation in dimensions ranging from 50 to 500.

\section{Background and Motivation}
Bayesian optimisation (BO) constructs a surrogate model---typically a Gaussian
process (GP)---of the objective function and uses an acquisition function to
select the next evaluation point~\cite{bibid}.  Standard BO methods assume
low-dimensional input spaces ($d \leq 20$) and fail in higher dimensions due to
the exponential growth of the search space and the poor conditioning of
covariance matrices~\cite{einstein1905}.

Recent advances in random projections, subspace identification, and trust-region
methods have extended BO to moderate dimensions ($20 < d \leq 100$) under the
assumption that the objective function varies primarily along a low-dimensional
manifold~\cite{goossensLaTeXCompanion1993}.  However, these methods rely on
rigid structural assumptions that do not hold for many real-world optimisation
problems in computational chemistry, materials design, and hyperparameter
tuning.

\section{Research Questions}
This doctoral research will address the following questions:

\begin{enumerate}
    \item \textbf{RQ1:} How can adaptive subspace identification be combined
    with trust-region methods to avoid premature convergence in high-dimensional
    Bayesian optimisation?
    \item \textbf{RQ2:} What information-theoretic acquisition functions best
    balance exploration and exploitation when the effective dimensionality of the
    objective function changes across the search space?
    \item \textbf{RQ3:} How do the proposed methods compare to existing
    high-dimensional BO approaches on benchmark functions and real-world
    engineering problems?
\end{enumerate}

\section{Methodology}

\subsection{Phase 1: Adaptive Subspace BO (Months 1--12)}
We will develop an adaptive subspace identification module that continuously
updates its estimate of the low-dimensional subspace as new function
evaluations become available.  The key innovation is a probabilistic subspace
prior that is updated via variational inference, enabling the model to track
shifts in the effective dimensionality of the objective function.

The acquisition function will be formulated in the embedded subspace:
\begin{equation}
    \alpha_{\text{sub}}(\mathbf{x}) = \mu_{\text{sub}}(\mathbf{Z}\mathbf{x})
    + \kappa \, \sigma_{\text{sub}}(\mathbf{Z}\mathbf{x}),
    \label{eq:ucb-sub}
\end{equation}
where $\mathbf{Z} \in \mathbb{R}^{k \times d}$ is the projection matrix
mapping from the original $d$-dimensional space to the $k$-dimensional subspace,
$\mu_{\text{sub}}$ and $\sigma_{\text{sub}}$ are the posterior mean and standard
deviation in the subspace, and $\kappa > 0$ controls the exploration--exploitation
trade-off.

\subsection{Phase 2: Information-Theoretic Extensions (Months 12--24)}
We will derive a family of information-theoretic acquisition functions
that directly maximise the expected reduction in entropy of the posterior
distribution over the global minimum.  These methods provide stronger
exploration guarantees than UCB-based approaches and are naturally adaptive to
changes in local dimensionality.

\subsection{Phase 3: Benchmarking and Applications (Months 24--36)}
The proposed methods will be evaluated on:
\begin{itemize}
    \item Synthetic benchmark functions (Branin, Hartmann,Ackley in high
    dimensions)
    \item Hyperparameter optimisation for deep learning models
    \item Molecular property optimisation in computational chemistry
    \item Aerodynamic shape optimisation using CFD simulations
\end{itemize}

\section{Timeline}
\begin{table}[htbp]
    \centering
    \caption{Research timeline and milestones.}
    \label{tab:timeline}
    \begin{tabular}{@{}lll@{}}
        \toprule
        \textbf{Period} & \textbf{Activity} & \textbf{Output} \\
        \midrule
        Months 1--12  & Adaptive subspace BO development & Journal paper \\
        Months 12--24 & Information-theoretic extensions  & Two conference papers \\
        Months 24--36 & Benchmarking and applications    & Final thesis chapters \\
        \bottomrule
    \end{tabular}
\end{table}

\section{Expected Contributions}
\begin{enumerate}
    \item A novel adaptive subspace BO algorithm with provable convergence
    guarantees in dimensions up to 500.
    \item A family of information-theoretic acquisition functions that adapt to
    varying local dimensionality.
    \item An open-source Python library implementing the proposed methods, with
    integration into the BoTorch framework.
    \item Comprehensive benchmarking results on standard and real-world
    optimisation problems.
\end{enumerate}

\section{Preliminary Work}
Preliminary experiments on the Hartmann-6 function demonstrate that the adaptive
subspace method reduces the simple regret by a factor of 3.2 compared to
standard high-dimensional BO after 200 function evaluations.

\printbibliography

\end{document}
